\section{Literature Review}

The literature relevant to this study sits at the intersection of Islamic finance, portfolio optimisation and machine learning-based financial prediction. 
From an Islamic finance perspective, investment decisions are not evaluated only through return maximisation, but also through their consistency with Shariah principles and broader ethical objectives. 
The Islamic Moral Economy framework argues that Islamic finance should move beyond formal legal compliance and contribute to social welfare, justice and human-centred development \citep{asutay2007conceptualisation, asutay2012conceptualising}. 
This is particularly important in portfolio construction, where the investment universe is first restricted through Shariah screening and then optimised according to risk-return objectives. 
In this regard, Shariah-compliant portfolio optimisation can be understood as a filtered investment process in which financial efficiency is pursued within ethical boundaries. 
While earlier studies on Islamic finance have emphasised the principles of risk sharing, prohibition of riba and asset-backed transactions \citep{elgamal2006islamic, ayub2009understanding, chong2009islamic}, 
recent developments in fintech and computational finance create new opportunities to operationalise these principles through data-driven investment models. 
Therefore, the use of machine learning and deep learning in Shariah-compliant portfolio management provides a practical bridge between the ethical foundations of Islamic finance and the technical requirements of modern financial markets.

